\subsection{持续分支熵率、projective ghost 全积分刚性与全息地址的 Lie 前缀失明}\label{subsec:conclusion-primeslice-ghost-holographic-lie-zerorate}
沿用定理 \ref{thm:conclusion-natural-extension-lstep-minimal-alphabet-submultiplicative}、定理 \ref{thm:conclusion-prime-adjunction-finite-depth-inverse-limit-budget-phase}、定理 \ref{thm:conclusion-primeslice-omega-support-chain-countable}、定理 \ref{thm:conclusion-primeslice-finite-depth-ledger-incompatibility}、定理 \ref{thm:conclusion-finite-padic-ghost-window-schur-rigidity}、定理 \ref{thm:conclusion-projective-ghost-family-recovers-valuation-histogram}、推论 \ref{cor:conclusion-modp-uniformity-forces-trivial-schur-sparsity}、定理 \ref{thm:conclusion-holographic-eventual-periodic-badic-rationals} 与定理 \ref{thm:conclusion-godel-filament-connected-lie-blindness} 的记号。前述结果已经分别给出：持续分支系统在任意有限深度上的 prime-slice 极小字母表与分支熵率、全深度 exact ledger 的 \(\omega\)-层最优性与有限秩加法账本障碍、有限 \(p^N\)-ghost 窗口对 \(Q_m\) 与全部 Schur 断层的模 \(p^N\) 刚性、projective ghost 塔对估值直方图与 Newton 多边形的无损恢复，以及全息地址线在 \(B\)-进有理点、周期骨架与连通 Lie 通道上的刚性分类。将这些对象交叉闭合后，还可进一步压缩为下列同一终端链。

\begin{theorem}[持续分支系统的 prime-slice 熵率]\label{thm:conclusion-branch-entropy-primeslice-rate}
\leanverified{paper\_conclusion\_branch\_entropy\_primeslice\_rate}
设 \(T:X\twoheadrightarrow X\) 为可数集合上的有限纤维满射，并记
\[
D_\ell(T):=\sup_{y\in X}\#T^{-\ell}(y)\in \NN\cup\{\infty\},
\qquad
h_{\mathrm{br}}(T):=\lim_{\ell\to\infty}\frac{1}{\ell}\log D_\ell(T)
=
\inf_{\ell\ge 1}\frac{1}{\ell}\log D_\ell(T).
\]
若对每个 \(\ell\ge 1\) 都有 \(D_\ell(T)<\infty\)，并定义长度-\(\ell\) 的极小单符号精确字母表位数为
\[
b_\ell^{\min}(T):=\Bigl\lceil \log_2 D_\ell(T)\Bigr\rceil,
\]
则有严格极限
\[
\lim_{\ell\to\infty}\frac{b_\ell^{\min}(T)}{\ell}
=
\frac{h_{\mathrm{br}}(T)}{\log 2}.
\]
特别地，若记相应极小字母表大小为
\[
\#\mathcal A_\ell^{\min}:=2^{b_\ell^{\min}(T)},
\]
则
\[
\#\mathcal A_\ell^{\min}
=
\exp\!\bigl(\ell h_{\mathrm{br}}(T)+o(\ell)\bigr).
\]
因此，当 \(h_{\mathrm{br}}(T)>0\) 时，任何长度-\(\ell\) 的精确单符号外置都必须支付指数级字母表。
\end{theorem}
\begin{proof}
由定理 \ref{thm:conclusion-natural-extension-lstep-minimal-alphabet-submultiplicative}，长度-\(\ell\) 的极小单符号精确字母表大小恰等于 \(D_\ell(T)\)，故
\[
b_\ell^{\min}(T)=\Bigl\lceil \log_2 D_\ell(T)\Bigr\rceil.
\]
同一定理还给出次乘法
\[
D_{\ell+r}(T)\le D_\ell(T)D_r(T),
\]
故 \((\log D_\ell(T))_{\ell\ge 1}\) 为次可加列，并由 Fekete 引理得到
\[
\lim_{\ell\to\infty}\frac1\ell \log D_\ell(T)=h_{\mathrm{br}}(T).
\]
另一方面，
\[
0\le b_\ell^{\min}(T)-\frac{\log D_\ell(T)}{\log 2}<1,
\]
故
\[
\frac{b_\ell^{\min}(T)}{\ell}
=
\frac{1}{\ell\log 2}\log D_\ell(T)+o(1)
\longrightarrow
\frac{h_{\mathrm{br}}(T)}{\log 2}.
\]
指数级字母表渐近式由
\[
2^{b_\ell^{\min}(T)}
=
\exp\!\bigl((\log 2)\,b_\ell^{\min}(T)\bigr)
\]
与上式立即得到。证毕。
\end{proof}

\begin{theorem}[\texorpdfstring{$\omega$}{omega}-层外置与有限秩线性化的双重阻断]\label{thm:conclusion-omega-support-finite-rank-double-obstruction}\leanverified{paper\_conclusion\_omega\_support\_finite\_rank\_double\_obstruction}
仍在定理 \ref{thm:conclusion-branch-entropy-primeslice-rate} 的设定下，若 \(T\) 在每一深度都存在真实分支，则成立以下两条互不替代的刚性：
\begin{enumerate}
  \item 任何精确编码全部无限前历史的 prime-slice 外置系统都必然需要 \(\omega\)-层有效 slices；不存在有限层 exact ledger。

  \item 即使在每个有限深度上精确外置已经实现，任意试图把这种精确性进一步压缩为某个有限生成交换群账本
  \[
  \Phi(n)=(e(n),s(n))\in \ZZ\oplus L
  \]
  的方案也必然失败。
\end{enumerate}
因此，持续分支系统的精确外置同时遭受两种互不替代的阻断：深度上必须无限分层，而代数上不能有限秩线性化。
\end{theorem}
\begin{proof}
第一条由定理 \ref{thm:conclusion-prime-adjunction-finite-depth-inverse-limit-budget-phase} 与定理 \ref{thm:conclusion-primeslice-omega-support-chain-countable} 直接给出：若每一深度都存在真实分支，则长度 \(N\) 前史的最小有效 prime-slice 数恰为 \(N\)，并且任意忠实全深度 prime-register 账本都必须使用可数无穷支撑，尤其不存在有限支持的精确外置。

第二条则是定理 \ref{thm:conclusion-primeslice-finite-depth-ledger-incompatibility} 的直接内容。该定理表明：有限深度 prime-slice 外置虽可严格精确，但任何试图把这种精确外置再压缩为某个有限生成交换群上的加法账本同态的方案都不可能成立。故有限秩线性化在代数层面被严格排除。证毕。
\end{proof}

\begin{theorem}[projective ghost 塔的全积分刚性]\label{thm:conclusion-projective-ghost-total-integral-rigidity}\leanverified{paper\_conclusion\_projective\_ghost\_total\_integral\_rigidity}
固定分辨率 \(m,m'\)。对每个素数 \(p\) 与每个 \(N\ge 1\)，记相应的 \(p^N\)-ghost 窗口为
\[
\mathcal G_{m,p^N},
\qquad
\mathcal G_{m',p^N}.
\]
若满足
\[
\mathcal G_{m,p^N}=\mathcal G_{m',p^N}
\qquad(\forall p,\ \forall N\ge 1),
\]
则成立以下全积分刚性：
\begin{enumerate}
  \item 对每个素数 \(p\) 与每个 \(r\ge 0\)，估值直方图一致：
  \[
  \mu_{m,p}(r)=\mu_{m',p}(r).
  \]

  \item 全部 \(p\)-进 Newton 多边形完全一致。

  \item 特征多项式在整数环中完全一致：
  \[
  Q_m(z)=Q_{m'}(z)\in \ZZ[z].
  \]

  \item 全部 Schur 断层作为整数表完全一致：
  \[
  T_{q,\lambda}(m)=T_{q,\lambda}(m')
  \qquad(\forall q\ge 0,\ \forall \lambda\vdash q).
  \]
\end{enumerate}
换言之，全 projective ghost 塔已经是 \(Q_m\) 及其全部 Schur 断层的完备整数不变量。
\end{theorem}
\begin{proof}
对固定的 \(p,N\)，定理 \ref{thm:conclusion-finite-padic-ghost-window-schur-rigidity} 给出
\[
Q_m(z)\equiv Q_{m'}(z)\pmod{p^N},
\qquad
T_{q,\lambda}(m)\equiv T_{q,\lambda}(m')\pmod{p^N}
\]
对每个 \(q,\lambda\) 都成立。
另一方面，定理 \ref{thm:conclusion-projective-ghost-family-recovers-valuation-histogram} 表明：整条 projective family \((\mathcal G_{m,p^N})_{N\ge 1}\) 无损恢复全部估值直方图 \((\mu_{m,p}(r))_{r\ge 0}\)，从而也无损恢复 \(Q_m\) 的完整 \(p\)-进 Newton 多边形。于是第 \(1\)、第 \(2\) 条立即成立。

对第 \(3\) 条，记
\[
Q_m(z)-Q_{m'}(z)=\sum_j c_j z^j.
\]
由上面的模 \(p^N\) 同余，对每个素数 \(p\) 与每个 \(N\ge 1\) 都有 \(p^N\mid c_j\)。因此每个整数 \(c_j\) 被任意高次 \(p\)-幂整除，只能为零。故 \(Q_m(z)=Q_{m'}(z)\) 于 \(\ZZ[z]\) 中成立。

对第 \(4\) 条，同理记
\[
T_{q,\lambda}(m)-T_{q,\lambda}(m')=:d_{q,\lambda}\in \ZZ.
\]
由定理 \ref{thm:conclusion-finite-padic-ghost-window-schur-rigidity}，对每个素数 \(p\) 与每个 \(N\ge 1\) 都有 \(p^N\mid d_{q,\lambda}\)。因此 \(d_{q,\lambda}=0\)。对所有 \(q,\lambda\) 逐一适用此论证，即得全部 Schur 断层整数表完全一致。证毕。
\end{proof}

\begin{theorem}[无限素数均匀影子的绝对 Schur 荒漠]\label{thm:conclusion-infinite-prime-uniform-shadow-absolute-schur-desert}
\leanverified{paper\_conclusion\_infinite\_prime\_uniform\_shadow\_absolute\_schur\_desert}
固定分辨率 \(m\)。若存在无穷多个奇素数 \(p\) 使得模 \(p\) 的非零残基直方图完全均匀，即
\[
N_a(m)=\nu_p
\qquad(\forall a\in \FF_p^\times),
\]
则对任意 \(q\ge 1\) 与任意分拆 \(\lambda\vdash q\)，都有
\[
T_{q,\lambda}(m)=0.
\]
因此，一旦“非零剩余类完全均匀”在无穷多素数上发生，则全部 Schur 断层整体塌缩为零。
\end{theorem}
\begin{proof}
固定 \(q\ge 1\) 与 \(\lambda\vdash q\)。由假设可取某个奇素数 \(p\) 满足
\[
p-1>q
\]
且模 \(p\) 的非零剩余类完全均匀。于是推论 \ref{cor:conclusion-modp-uniformity-forces-trivial-schur-sparsity} 适用，并给出
\[
T_{q,\lambda}(m)\equiv 0\pmod p.
\]
由于满足这些条件的素数 \(p\) 无穷多，而 \(T_{q,\lambda}(m)\) 是固定整数，故它只能为零。对任意 \(q,\lambda\) 均如此，命题遂证。证毕。
\end{proof}

\begin{theorem}[全息地址线的周期骨架稠密而零测]\label{thm:conclusion-holographic-periodic-skeleton-dense-null}
沿用小节 \ref{subsec:conclusion-holographic-localized-zg-rigidity} 的记号，定义
\[
X_E:=\mathcal X_E\subset T_v\cong \ZZ_2,
\qquad
\mathrm{Per}_B:=\QQ_B v\cap X_E.
\]
对任意满足
\[
H(\mathbf p)>0
\]
的 Bernoulli 源 \(\mathbf p=(p_a)_{a\in E}\)，记其全息推前测度为 \(\mu_{\mathbf p}\)。则集合 \(\mathrm{Per}_B\) 具有双重性质：
\begin{enumerate}
  \item \(\mathrm{Per}_B\) 在 \(X_E\) 中稠密；
  \item \(\mu_{\mathbf p}(\mathrm{Per}_B)=0\)。
\end{enumerate}
因此，window-free 全息地址线同时具有一个算术上稠密而测度上稀薄的周期骨架。
\end{theorem}
\begin{proof}
由定理 \ref{thm:conclusion-holographic-eventual-periodic-badic-rationals}，
\[
x\in X_E\ \text{来自某条最终周期事件流}
\iff
x\in \QQ_Bv.
\]
因此 \(\mathrm{Per}_B\) 正是最终周期事件流在全息地址映射下的像。

再由定理 \ref{thm:xi-time-part62d-hologram-image-cantor-homeomorphism}，地址映射
\[
X:E^{\NN}\longrightarrow X_E
\]
是拓扑同胚。由于最终周期序列在移位空间 \(E^{\NN}\) 中稠密，故其像 \(\mathrm{Per}_B\) 在 \(X_E\) 中也稠密。

另一方面，最终周期序列集是可数的，故 \(\mathrm{Per}_B\) 亦可数。若 \(H(\mathbf p)>0\)，则 Bernoulli 源非退化，因而乘积测度在 \(E^{\NN}\) 上对每个单点都赋零质量。由 \(X\) 的双射性可知，\(\mu_{\mathbf p}\) 对 \(X_E\) 的每个单点同样赋零质量。可数并集仍为零测，于是
\[
\mu_{\mathbf p}(\mathrm{Per}_B)=0.
\]
证毕。
\end{proof}

\begin{theorem}[连通 Lie 通道的有限前缀因子化与零信息率]\label{thm:conclusion-connected-lie-prefix-factorization-zero-rate}
设
\[
f:T_v\longrightarrow G
\]
为到任意有限维连通实 Lie 群 \(G\) 的连续群同态。则存在整数 \(N\ge 0\)、有限 \(2\)-群 \(Q\) 与因子化
\[
T_v\twoheadrightarrow T_v/2^N T_v\xrightarrow{\ \bar f\ } Q\hookrightarrow G.
\]
并且存在某个前缀深度 \(n_0\ge 0\) 使得对任意两条事件流 \(\mathbf e,\mathbf e'\in E^{\NN}\)，若其前 \(n_0\) 个编码数字相同，则
\[
f(X(\mathbf e))=f(X(\mathbf e')).
\]

更进一步，若 \(\mathbf e\sim \nu_{\mathbf p}\) 为非退化 Bernoulli 源，并记
\[
Y:=f(X(\mathbf e)),
\]
则
\[
\sup_{n\ge 1} I\bigl((e_1,\dots,e_n);Y\bigr)\le \log|Q|<\infty,
\]
从而
\[
\lim_{n\to\infty}\frac1n I\bigl((e_1,\dots,e_n);Y\bigr)=0.
\]
因此，任何连通 Lie 相位通道都只能读取有限前缀，且其每符号信息率必为零。
\end{theorem}
\begin{proof}
由定理 \ref{thm:conclusion-godel-filament-connected-lie-blindness}，像群
\[
Q:=f(T_v)
\]
必为有限 \(2\)-群。又因为 \(T_v\cong \ZZ_2\)，其开子群恰为 \(2^N T_v\) 的形式。定理 \ref{thm:conclusion-godel-filament-connected-lie-blindness} 进一步给出 \(\ker f\) 为开子群，故存在 \(N\ge 0\) 使
\[
\ker f=2^N T_v.
\]
于是 \(f\) 必经由有限商 \(T_v/2^N T_v\) 因子化为
\[
T_v\twoheadrightarrow T_v/2^N T_v\xrightarrow{\ \bar f\ } Q\hookrightarrow G.
\]

同一定理还说明：存在 \(n_0\ge 0\) 使得只要两条事件流前 \(n_0\) 个编码数字相同，其地址差便落在 \(\ker f\) 中，从而在 \(f\circ X\) 下具有相同像。故 \(Y\) 实际上只是有限前缀 \((e_1,\dots,e_{n_0})\) 的函数。

因此对任意 \(n\ge 1\)，有
\[
I\bigl((e_1,\dots,e_n);Y\bigr)
=
I\bigl((e_1,\dots,e_{\min(n,n_0)});Y\bigr)
\le
H(Y).
\]
而 \(Y\) 的值域包含在有限群 \(Q\) 中，故
\[
H(Y)\le \log|Q|.
\]
由此得到
\[
\sup_{n\ge 1}I\bigl((e_1,\dots,e_n);Y\bigr)\le \log|Q|<\infty.
\]
再两边除以 \(n\) 并令 \(n\to\infty\)，便得零信息率结论。证毕。
\end{proof}

\begin{conjecture}[正信息率可观测量必须离开连通 Lie 世界]\label{conj:conclusion-positive-rate-observable-leaves-connected-lie}
设 \(Y=F(X(\mathbf e))\) 是定义在全息地址线 \(X_E\) 上的连续群值观测，并在某个非退化 Bernoulli 源下满足
\[
\limsup_{n\to\infty}\frac1n I\bigl((e_1,\dots,e_n);Y\bigr)>0.
\]
则 \(Y\) 的像所生成的拓扑群不可能是有限维连通 Lie 群；更强地说，它必须保留某个无穷全不连通的 \(2\)-进因子。

这一判断可视为定理 \ref{thm:conclusion-connected-lie-prefix-factorization-zero-rate} 的逆向强化：连通 Lie 通道已经被证明必为零信息率，因此一切真正携带正熵率的可观测量都应被强制推出连通 Lie 包络，而落入 profinite 或 \(2\)-adic 方向。该图景与定理 \ref{thm:conclusion-localization-visible-circle-blindness-slicegrowth} 所揭示的“连通 Lie 可见量对 slice growth 的系统性失明”完全一致。
\end{conjecture}

上述链条可压缩为
\[
\boxed{
\text{持续分支熵率 }h_{\mathrm{br}}(T)
\Longrightarrow
\text{prime-slice 字母表指数律}
\Longrightarrow
\omega\text{-层 exact ledger 不可有限化}
\Longrightarrow
\text{projective ghost 塔的全积分刚性}
}
\]
\[
\boxed{
\text{无限素数均匀影子}
\Longrightarrow
\text{absolute Schur desert}
\Longrightarrow
\text{全息地址线的稠密零测周期骨架}
\Longrightarrow
\text{连通 Lie 通道的有限前缀化与零信息率}
}
\]
因此，精确外置成本、\(p\)-进 ghost 刚性与连通 Lie 可见性的坍缩并非三条并列理论，而是同一结构事实的三种投影：可交换可见层可以极度贫瘠，而全不连通算术层却同时承载指数级历史、完整 \(p\)-进谱表与稠密周期骨架。
